\Needspace{18\baselineskip}
\subsection{Models pulled back by a finite curve map}
\label{subsec:finite-curve-pullback-models}

Let $r\colon C'\to C$ be a finite surjective holomorphic map of
degree $d$ between connected compact Riemann surfaces. Fix a
smooth real form $\beta$ of area $A\geq0$ on $C$ and put
$\sigma=r^*\beta$. The definitions and results of this section
apply on $C$ and on $C'$ in place of $Y$. The models pulled back
from $C$ occupy a proper subspace of the full model space on $C'$
in general.
The fiber maximum of \cref{finite-curve:fiber-maximum} compares
envelope competitors on the two curves. The weighted measure
pullback then identifies the metric on this subspace and gives
a retraction from the full space.

\Needspace{15\baselineskip}
\begin{proposition}\label{finite-curve:envelope-pullback}\index{Envelope!finite curve pullback}
For $v\in\PSH(C,\beta)$ and every smooth real obstacle $F$ on $C$,
\[
 \begin{aligned}
 P_{\sigma,F\circ r}[v\circ r]
       &=P_{\beta,F}[v]\circ r,\\
 P_{\sigma,F\circ r}[v\circ r]_{\mathcal I}
       &=P_{\beta,F}[v]_{\mathcal I}\circ r.
 \end{aligned}
\]
In particular, $J(v)=v\circ r$ embeds the ordinary model space
$\mathcal M(\beta)$ into $\mathcal M(\sigma)$ and preserves
ideal projection.
\end{proposition}
\begin{proof}
Pullback sends every base competitor to an upstairs competitor:
it preserves positivity, the obstacle inequality and the ordinary
order. It also preserves and reflects the all-weight ideal order,
because on a curve this order consists of the pointwise Lelong
inequalities, as in the proof of \cref{finite-curve:fiber-maximum},
and $\nu(w\circ r,x)=e_x\nu(w,r(x))$ by
\cref{curve-map:lelong-pullback}, where $e_x$ is the local degree
of $r$ at $x$.
Conversely, let $h$ be an upstairs competitor, with fiber maximum
$m_h$. Then $h\leq m_h\circ r$, and $h\leq F\circ r$ implies
$m_h\leq F$. By \cref{finite-curve:fiber-maximum},
$h\preceq v\circ r$ implies $m_h\preceq v$, and the same holds
for the ideal order. For an ideal competitor, moreover,
$v\circ r\preceq_{\mathcal I}h\preceq_{\mathcal I}m_h\circ r$,
and reflection gives $v\preceq_{\mathcal I}m_h$. Thus $m_h$ is a
base competitor in both cases, and the two unregularized suprema
are equal after pullback.
A finite nonconstant holomorphic map of curves is open;
upper regularization commutes with its pullback, as is seen
in the coordinates $t=z^e$. This proves both identities.
Take $F=0$. An ordinary model is fixed by its ordinary
envelope, so its pullback is a model; surjectivity of $r$
gives injectivity. Applying the ideal identity to a model
proves preservation of ideal projection.
\end{proof}

Let $d_\beta$ and $d_\sigma$ be the polar metrics from
\cref{prop:curve-arbitrary-background-models}. If $q$ is a model
on $C'$, denote its polar curvature by
$\mu_q=\mathbf1_{\{q=-\infty\}}(\sigma+\ddc q)$.
Write $p_\beta(\nu)$ for the base model with polar measure
$\nu$, using the same classification.
Write $U_\beta(v)=\int_C\mathrm{MA}_\beta(v)$ and
$U_\sigma(q)=\int_{C'}\mathrm{MA}_\sigma(q)$ for their
non-pluripolar volumes.

\Needspace{18\baselineskip}
\begin{proposition}\label{finite-curve:model-retraction}\index{Metric completeness!finite curve pullback subspace}\index{Model potential!finite curve averaging}
For $u,v\in\mathcal M(\beta)$,
\[
 d_\sigma(J(u),J(v))=d\,d_\beta(u,v),\qquad
 U_\sigma(J(u))=d\,U_\beta(u).
\]
The image $J(\mathcal M(\beta))$ is closed and contains the
linear polar-measure geodesic between each pair of its members.
The map defined for $q\in\mathcal M(\sigma)$ by
\[
 R(q)=J\!\left(p_\beta\!\left(d^{-1}r_*\mu_q\right)\right)
\]
is a $1$-Lipschitz retraction onto this image. Its polar measure is
\[
 \mu_{R(q)}=d^{-1}r^*r_*\mu_q.
\]
It preserves the polar mass and the non-pluripolar volume,
and commutes with ideal projection on $\mathcal M(\sigma)$.
At $A=0$ both model spaces are smooth singletons.
\end{proposition}
\begin{proof}
Change of variables off the finite branch sets gives
$\int_{C'}\sigma=dA$.
By \cref{finite-curve:measure-pullback}, the polar measure of
$J(v)$ is $r^*\mu_v$, with total mass $d\mu_v(C)$.
The curve classification gives
\[
 \begin{aligned}
 d_\sigma(J(u),J(v))
   &=\|r^*(\mu_u-\mu_v)\|_{\mathrm{TV}}
     =d\|\mu_u-\mu_v\|_{\mathrm{TV}},\\
 U_\sigma(J(v))
   &=dA-d\mu_v(C)=d\,U_\beta(v).
 \end{aligned}
\]
Completeness of the base model space proves closedness of the
image: a convergent sequence in the image comes from a
Cauchy sequence downstairs. Linear interpolation of polar
measures commutes with $r^*$, giving the stated geodesics.

For any $q\in\mathcal M(\sigma)$, its polar measure has mass
at most $dA$. Its pushforward is carried by the image of a polar
carrier, which is polar by the image assertion
of \cref{finite-curve:measure-pullback}, so
$d^{-1}r_*\mu_q$ has mass at most $A$ and defines a base model.
This proves that $R$ is well defined and gives its displayed
polar measure. The identity $r_*r^*=d$ makes $R\circ J=J$.
For any finite signed measure $\tau$ on $C'$,
pushforward does not increase total variation, hence
\[
 \|d^{-1}r^*r_*\tau\|_{\mathrm{TV}}
       =\|r_*\tau\|_{\mathrm{TV}}
       \leq\|\tau\|_{\mathrm{TV}}.
\]
Apply this to the difference of two polar measures to obtain
the Lipschitz estimate.
The averaging operator $d^{-1}r^*r_*$ preserves the total mass
of positive measures, so the formula $U_\sigma(q)=dA-\mu_q(C')$
gives the preservation of polar mass and volume.
Both measure operations commute with taking the atomic part,
which is ideal projection by the curve classification.
This proves the commutation assertion. At zero area all
polar measures vanish and both spaces are singletons.
\end{proof}

This retraction averages the singularity data over a fiber.
Unlike the ideal retraction, it need not select a nearest model.

\Needspace{15\baselineskip}
\begin{example}\label{finite-curve:averaging-not-nearest}\index{Model potential!averaging is not nearest}
For $r\colon\mathbb P^1\to\mathbb P^1$, $r(z)=z^3$, and the
Fubini--Study background $\beta$ of area one, the retraction
$R$ above is not a nearest-point retraction, even on ideal
models with analytic singularities.
\end{example}
\begin{proof}
Put $\sigma=r^*\beta$, of area three, and let $q$ be its model
with polar measure $\delta_1$. This is an analytic ideal model:
in the affine coordinate, the global function
\[
 h(z)=\log|z-1|^2-\tfrac13\log(1+|z|^6)
\]
extends smoothly across infinity and satisfies
$\sigma+\ddc h=\delta_1+\tfrac23\sigma$.
Thus $h$ is $\sigma$-plurisubharmonic with one logarithmic pole,
and the curve classification identifies its envelope with $q$;
\cref{prop:analytic-envelope-type} preserves its analytic type.
Writing $\zeta=\exp(2\pi\mathrm i/3)$, one has
\[
 \mu_{R(q)}=\tfrac13(\delta_1+\delta_\zeta+\delta_{\zeta^2}).
\]
The three atoms are distinct, so
\[
 d_\sigma(q,R(q))=\tfrac23+\tfrac13+\tfrac13=\tfrac43.
\]
The zero model is $J(0)$ and has zero polar measure, giving
$d_\sigma(q,J(0))=1$. Hence the averaged model is strictly
farther from $q$ than a member of the same pullback subspace.
\end{proof}
